\section{Corpus Self-Contradictions}
\label{app:contradictions}
\setlength{\emergencystretch}{2em}

This appendix expands \S\ref{sec:discussion} (conflicting guidance) with three of the 16 confirmed corpus self-contradictions. For each query we show the \emph{two conflicting passages} the model retrieved, each in its own block labelled with its (openly available) source and page, the conflicting statement highlighted. Both reach the generator in a single context, so whichever it follows can be faulted against the other.

\subsection*{Re-dosing emergency contraception after vomiting}
\noindent\textbf{Query.} \emph{How should high-dose emergency contraceptive pills be dosed and managed if vomiting occurs?}\quad{\scriptsize\ttfamily q\_00459}
\smallskip
\noindent\textbf{Hesperian, A Book for Midwives (p.318)}
\begin{quote}\footnotesize
\textbf{Using high-dose pills (with 50 mcg ethinyl estradiol)}\\
\textbullet\ give 2 pills. Then 12 hours later, give 2 more pills.\\
Emergency pills can give women headaches or severe nausea. Women can try eating something at the same time as taking the pills. \hl{If a woman vomits within 3 hours of taking the pills, she should take the same dose again.}
\end{quote}
\noindent\textbf{WHO contraception guidance (p.74)}
\begin{quote}\footnotesize
\textbf{3. Explain what to do about side effects}\\
\textbullet\ Nausea:\\
\textbullet\ Routine use of anti-nausea medications is not recommended.\\
\textbullet\ Women who have had nausea with previous ECP use or with the first dose of a 2-dose regimen can take anti-nausea medication such as 25–50 mg meclizine hydrochloride (such as Agyrax, Antivert, Bonine, Postafene) one-half to one hour before taking ECPs.\\
\textbullet\ Vomiting:\\
\textbullet\ \hl{If the woman vomits within 2 hours after taking progestin-only or combined ECPs, she should take another dose.} If she vomits within 3 hours of taking ulipristal acetate ECPs, she should take another dose. (She can use antinausea medication with this repeat dose, as above.) If vomiting continues, she can take a repeat dose of progestin-only or combined ECPs by placing the pills high in her vagina. If vomiting occurs more than 2 hours after taking progestin-only or combined ECPs, or 3 hours after taking UPA-ECPs, then she does not need to take any extra pills.
\end{quote}
\noindent\textbf{The conflict.} A re-doses if vomiting occurs within \textbf{3 hours}; B (for the same progestin-only or combined pills) within \textbf{2 hours}.

\subsection*{The cervical dilatation that starts active labour}
\noindent\textbf{Query.} \emph{At what cervical dilation should partograph plotting commence?}\quad{\scriptsize\ttfamily q\_00560}
\smallskip
\noindent\textbf{WHO Midwifery Education Modules (p.65)}
\begin{quote}\footnotesize
\textbf{Progress of labour}\\
Cervical dilatation The first stage of labour is divided into the latent and active phases.\\
The latent phase at the onset of labour lasts until cervical dilatation is 4 cm and is accompanied by effacement of the cervix. The latent phase may last up to 8 hours, although it is usually completed more quickly than this. Although regular assessments of maternal and fetal well-being and a record of all findings should be made, these are not plotted on the partograph (using modified version) until labour enters active phase.\\
\hl{The active phase of the first stage of labour starts when the cervix is 4 cm dilated} and is completed at full dilatation, i.e. 10 cm. Progress in this phase is approximately 1 cm per hour and often quicker in multigravidae.
\end{quote}
\noindent\textbf{WHO Labour Care Guide (p.11)}
\begin{quote}\footnotesize
\textbf{When should the LCG be initiated?}\\
\hl{When women have entered the active phase of the first stage of labour (i.e. cervical dilatation of 5 cm or more).}
\end{quote}
\noindent\textbf{The conflict.} A puts the active phase (and partograph plotting) at \textbf{4\,cm}; B at \textbf{5\,cm or more}.

\subsection*{How long after birth a placenta is ``retained''}
\noindent\textbf{Query.} \emph{How long after birth is a placenta considered retained?}\quad{\scriptsize\ttfamily q\_00761}
\smallskip
\noindent\textbf{WHO Midwifery Education Modules (p.78)}
\begin{quote}\footnotesize
\textbf{haemorrhage bleeding.}\\
If the bleeding is atonic it is important to find out if the placenta has been delivered or not. When the third stage is managed actively, the placenta is normally delivered within 5–10 minutes of birth of the baby. It takes longer to separate and deliver during physiological management, 20–30 minutes.\\
\hl{The students should remember that the placenta is retained if it has not been delivered within one hour of the delivery of the baby.}\\
If the placenta has not been delivered, is there bleeding?
\end{quote}
\noindent\textbf{NICE intrapartum guideline (p.69)}
\begin{quote}\footnotesize
\textbf{Prolonged third stage}\\
1.10.20 \hl{Diagnose a prolonged third stage of labour if it is not completed within 30 minutes of the birth with active management or within 60 minutes of the birth with physiological management.} Follow the recommendations on managing a retained placenta. [2014]
\end{quote}
\noindent\textbf{The conflict.} A calls a placenta retained after \textbf{1 hour} regardless of management; B uses \textbf{30 minutes} (active) / \textbf{60 minutes} (physiological).
